% Discussion (Task 42): handoff/37_narrative_v2.md section 3.4 applied; implications paragraph added (plan 32).
\section{Discussion}
\label{sec:discussion}

\paragraph{Stable but not faithful.}
On both datasets, rollout attention agreed weakly to moderately with attribution methods and kept most of its
top-\paramTopK{} mass under the feature constraints, yet it was clearly less faithful than IG and occlusion by deletion
and insertion AUC (Section~\ref{sec:rq1}).
An analyst who reads rollout attention as feature importance on these models would therefore see a stable picture that
only partly reflects what drives the prediction.
This holds for the two datasets and the FT-Transformer we studied; we do not claim that attention is unfaithful in
general.

\paragraph{Stiffness is not a security property.}
The small change of attention under A2 and A3 is a statement about the feature constraints and budgets we used.
Among the \paramTopK{} rollout features with the highest mass on clean explanation rows, on average
\topImmRolloutSeventeen{} on 2017 and \topImmRolloutEighteen{} on 2018 were immutable, and among the clean
top-\paramTopK{} IG features that A2 targets, \topImmIGSeventeen{} and \topImmIGEighteen{} (macro over the
\nCommonClasses{} classes, mean over seeds).
The stiffness is therefore partly structural, because the attacker cannot change these features.
The unconstrained reference moved attention further, so stronger or differently constrained attackers might move it
more.

\paragraph{The faithfulness loss follows the prediction flip.}
The faithfulness loss after an evasion attack appeared for most explanation methods, including the static class-mean
rankings, which do not depend on the input.
A plausible reading is that adversarial rows sit in a region where the model's output depends on the features
differently from clean rows, so any fixed explanation becomes less faithful there.
The comparison of flipped rows, unflipped rows and rows flipped by the random control supports this reading; the
mechanism itself was not tested.

\paragraph{Consistency monitoring.}
Consistency scores helped against an attacker that ignored them; against an attacker that optimised the detector's
exact score they did not help, and D$_+$ was easier to evade than D (Section~\ref{sec:rq2}).
Against the attacker that ignored them, the gain was not specific to attention: controls that used other attribution
methods instead of attention, or statistics of the logits and distances, obtained most or all of it, and a1 still
carries the class information of the true label.
The small gain that D$_+$ showed against A8 did not hold against A8x, which uses exact IG in the loop and selects
iterates by the true score: the sign reversed on 2017 in the pooled estimate and in all but \axHeadOppSeedsSeventeen{} of the \nSeeds{} seeds,
and giving D's attacker the same
selection changed no row.
An evaluation with stand-ins alone would therefore have overstated what D$_+$ achieves; this is the general point made by
\cite{tramer2020}, shown here for consistency monitoring.
We did not test why D$_+$ was easier to evade than D.
One untested reading is that each added score gives the attacker a further direction to move in while the threshold
stays at \paramFPR{} FPR; the label-shuffled control a1, under its own stand-in attacker, also left more rows
undetected than D on 2017 (\cmdAoneVoneSeventeen, \cmdAoneVoneSeventeenCI), but a1 uses the true class label and a
weaker attacker, so it does not test this reading.

\paragraph{Dataset dependence.}
The difference between the datasets in detector strength was much larger than the difference between D and D$_+$.
On 2017 a single class dominated the successful and undetected rows under A6, and on 2018 a different single class did
(Section~\ref{sec:rq3}).
Conclusions about learned attack detectors on these benchmarks should therefore be reported per dataset and per class.

\paragraph{Implications for NIDS practice.}
For an analyst, attention maps of an FT-Transformer flow classifier are a stable but partial picture of what drives a
prediction; when a decision rests on feature importance, IG or occlusion were the more faithful choice on both
datasets, and any explanation of a row that an attack has flipped is less faithful than on clean traffic.
For a defender, we do not recommend these consistency scores as an addition to a confidence-and-distance detector
against adaptive attackers: they helped only against an attacker that ignored them, that help was available without
attention, and against an attacker that modelled them exactly they made evasion easier on both datasets (on 2018
close to the floor).
For anyone evaluating an explanation-based attack detector, the attacker should optimise the detector's exact score
wherever a gradient exists and select its iterates by the true score; an evaluation with differentiable stand-ins can
reverse the sign of the result, as it did here.
